%%%%%%%%%%%%%%%%%%%%%%%%%%%%%%%%%%%%%%%%%%%%%%%%%%%%%%%%%%%%%%%%%%%%%%%%%%%%%%%
%
%  customcommands.tex -- Zusatzpakete und eigene Befehle
%
%  Arbeitsteilung der Vorlage:
%    structure_document.tex  Seitenlayout, Schriften, Ueberschriften, Boxen
%    customcommands.tex      dieses File: Komfortpakete und eigene Makros
%    acronyms.tex            Glossar- und Abkuerzungseintraege
%
%  ACHTUNG: Diese Datei wird NACH structure_document.tex eingebunden und
%  setzt Definitionen von dort voraus -- insbesondere \avantfont (Schrift der
%  Titelseite) und die Farbe ThemeColor.
%
%%%%%%%%%%%%%%%%%%%%%%%%%%%%%%%%%%%%%%%%%%%%%%%%%%%%%%%%%%%%%%%%%%%%%%%%%%%%%%%


%==============================================================================
%  1. PAKETE
%==============================================================================

%--- Grafiken -----------------------------------------------------------------
\usepackage{graphicx}
\graphicspath{{../Pictures/}} % alle Bilder liegen hier; Angabe im Text ohne "Pictures/"
\usepackage{tikz}          % Titelseite, Kapitelkoepfe, Boxen
\usepackage{etoolbox}      % robuste Tests wie \ifdefempty (s. Titelseite)

%--- Listen und Tabellen ------------------------------------------------------
\usepackage{enumitem}      % Listen anpassbar
\setlist{nolistsep}        % kompaktere Abstaende in itemize/enumerate
\usepackage{booktabs}      % \toprule, \midrule, \bottomrule statt \hline

%--- Gleitobjekte -------------------------------------------------------------
% Zwingt Abbildungen und Tabellen, vor dem naechsten \section zu erscheinen.
% Wirkt nicht auf \subsection -- dort bei Bedarf \FloatBarrier von Hand setzen.
\usepackage[section]{placeins}

%--- Nur fuer die Beispielkapitel ---------------------------------------------
% \lipsum erzeugt Blindtext. Fuer die Abgabe-/Release-Version kann diese Zeile zusammen mit Document/99_example.tex entfernt werden.
\usepackage{lipsum}

%==============================================================================
%  2. ABBILDUNGEN
%==============================================================================

%------------------------------------------------------------------------------
%  \ofig -- "One FIGure": Abbildung mit einem Bild
%------------------------------------------------------------------------------
%  Nutzung:
%    \ofig{Bildunterschrift}{dateiname}{breite}
%    \ofig[eigenesLabel]{Bildunterschrift}{dateiname}{breite}
%
%  Verweis darauf:  \ref{fig:dateiname}  bzw.  \ref{fig:eigenesLabel}
%
%  Der Dateiname wird ohne Endung und ohne "Pictures/" angegeben.
%  Breite z.B. \linewidth, 0.6\linewidth, 8cm.
%
%  Ohne optionales Label wird der Dateiname als Label benutzt. Das ist bequem,
%  geht aber schief, sobald der Dateiname Unterstriche oder Punkte enthaelt --
%  dann das optionale Argument verwenden.
%------------------------------------------------------------------------------
\NewDocumentCommand{\ofig}{o m m m}{%
	\begin{figure}[htbp]
		\centering
		\includegraphics[width=#4]{#3}%
		\caption{#2}%
		\IfValueTF{#1}{\label{fig:#1}}{\label{fig:#3}}%
	\end{figure}%
}

%------------------------------------------------------------------------------
%  \pfig -- "Placeholder FIGure": Abbildung mit dem Platzhalterbild
%------------------------------------------------------------------------------
%  Nutzung:  \pfig{Bildunterschrift}
%  Praktisch, solange das echte Bild noch fehlt, damit Umbruch und
%  Nummerierung schon stimmen.
%------------------------------------------------------------------------------
\NewDocumentCommand{\pfig}{o m}{%
	\IfValueTF{#1}{\ofig[#1]{#2}{placeholder}{\linewidth}}%
	              {\ofig{#2}{placeholder}{\linewidth}}%
}


%==============================================================================
%  3. TITELSEITE
%==============================================================================

%------------------------------------------------------------------------------
%  Stellschrauben der Titelseite
%------------------------------------------------------------------------------
%  Diese drei Werte steuern das Layout. Die Bildhoehe wird daraus berechnet,
%  damit die Titelseite auch bei anderem Papierformat als A4 stimmt:
%  Die Oberkante des Bildes liegt bei 3/5 der Seitenhoehe, der farbige Balken
%  steht darunter am Seitenfuss.
%------------------------------------------------------------------------------
\newcommand{\hkaCoverImage}{cover_image.jpg} % leer lassen => kein Titelbild

\newlength{\hkaColorBlockHeight}
\setlength{\hkaColorBlockHeight}{3cm}        % Hoehe des farbigen Balkens unten

\newlength{\hkaCoverImageHeight}
\setlength{\hkaCoverImageHeight}{\dimexpr 0.6\paperheight - \hkaColorBlockHeight\relax}
% Bei A4: 0.6 * 29.7cm = 17.82cm, minus 3cm Balken = 14.82cm.

%------------------------------------------------------------------------------
%  \createTikZTitlePage -- baut die Titelseite auf
%------------------------------------------------------------------------------
%  Nutzung:
%    \createTikZTitlePage{Fakultaet}{Titel}{Autor(en)}
%
%  Der Fakultaetsname darf mit \\ umgebrochen werden.
%  Die Farbe kommt aus ThemeColor, wird also in document.tex festgelegt.
%------------------------------------------------------------------------------
\NewDocumentCommand{\createTikZTitlePage}{m m m}{%
	\begingroup
	\thispagestyle{empty} % keine Kopf-/Fusszeile auf der Titelseite
	\begin{tikzpicture}[remember picture,overlay]

	%--- Ankerpunkte (absolut auf der Seite) --------------------------------
		\coordinate (UoAS)         at ($ (current page.north west) + (2,-2) $);
		\coordinate (Faculty)      at ($ (UoAS) + (0,-2) $);
		\coordinate (Title)        at ($ (current page.north west) + (2,-11.5) $);
		\coordinate (Authors)      at ($ (current page.south west) + (2,1.5) $);
		\coordinate (BackgroundPic) at ($ (current page.south) + (0,\hkaColorBlockHeight) $);
		% Ankerpunkte des "+|KA"-Logos
		\coordinate (H) at ($ (current page.north east) + (8,-0.8) $);
		\coordinate (K) at ($ (H) + (0,-18) $);
		\coordinate (A) at ($ (K) + (0,-8.6) $);

	%--- Titelbild ----------------------------------------------------------
	% Ist \hkaCoverImage leer, bleibt die Flaeche frei -- so laesst sich die
	% Vorlage auch ohne mitgeliefertes Bild uebersetzen.
		\ifdefempty{\hkaCoverImage}{}{%
			\node[above, inner sep=0pt] (Hinterlegbild) at (BackgroundPic)
				{\includegraphics[height=\hkaCoverImageHeight]{\hkaCoverImage}};%
		}

	%--- Farbiger Balken am Seitenfuss --------------------------------------
		\path[fill=ThemeColor] (current page.south west)
			rectangle ($ (current page.south east) + (0,\hkaColorBlockHeight) $);

	%--- Textblöcke ---------------------------------------------------------
		% Hochschule
		\node[text width = 10cm, ThemeColor, below right] at (UoAS)
			{\avantfont{\textbf{Hochschule Karlsruhe}}\\
			 \avantfont{University of}\\
			 \avantfont{Applied Sciences}};
		% Fakultät
		\node[text width = 10cm, ThemeColor, below right] at (Faculty)
			{\avantfont{Fakultät für}\\ \avantfont{\textbf{#1}}};
		% Titel (und ggf. Untertitel)
		\node[text width = 10cm, ThemeColor, above right] at (Title) {#2};
		% Autor(en) -- steht im farbigen Balken, daher weiß
		\node[text width = 10cm, white, right] at (Authors) {{\large #3}};

	%--- Logo "+|KA" --------------------------------------------------------
	% Als TikZ-Pfade gezeichnet, damit das Logo scharf und in ThemeColor
	% bleibt. "+(x,y)" ist relativ zum zuletzt gesetzten Ankerpunkt, ohne
	% diesen zu verschieben.
		\path (H) [fill=ThemeColor, scale = 0.45, rounded corners = 1]
			+(0,-2.95) -- +(3,-2.95) -- +(3,0) -- +(4.2,0) -- +(4.2,-2.95)
			-- +(7.2,-2.95) -- +(7.2,-3.95) -- +(4.2,-3.95) -- +(4.2,-6.9)
			-- +(3,-6.9) -- +(3,-3.95) -- +(0,-3.95) -- cycle;
		\path (H) [fill=ThemeColor, scale = 0.45, rounded corners = 1]
			+(7.8,0) -- +(9,0) -- +(9,-6.9) -- +(7.8,-6.9) -- cycle;
		\path (K) [fill=ThemeColor, scale = 0.45, rounded corners = 1]
			+(3,0) -- +(4.2,0) -- +(4.2,-2.8) -- +(9,0) -- +(9,-1.3)
			-- +(5.2,-3.45) -- +(9,-5.6) -- +(9,-6.9) -- +(4.2,-4.1)
			-- +(4.2,-6.9) -- +(3,-6.9) -- cycle;
		\path (A) [fill=ThemeColor, scale = 0.45, rounded corners = 1]
			+(6,0) -- +(6.65,0) -- +(10.05,-6.9) -- +(8.75,-6.9) -- +(6,-1.3)
			-- +(4.8,-3.7) -- +(7.52,-3.7) -- +(8.85,-4.8) -- +(4.3,-4.8)
			-- +(3.25,-6.9) -- +(1.95,-6.9) -- +(5.35,0) -- cycle;

	\end{tikzpicture}%
	\vfill
	\endgroup
}